\chapter{Imparare alla presenza di qualcuno}\label{ch:synthesis}

Il Portale prende il proprio nome sul serio quando il vetro smette di essere soltanto una superficie per contenuti e diventa il luogo in cui un'altra persona può apparire, ascoltare e rispondere. La scoperta non consiste nel dimenticare che ci separa una tecnologia. Consiste nell'accorgersi che, attraverso quella tecnologia, possiamo fare qualcosa insieme che prima richiedeva di essere nello stesso luogo.

L'immagine quasi speculare dà una forma a questa possibilità. Due persone occupano posizioni corrispondenti senza diventare uguali. Una conosce un passaggio, l'altra propone un ritmo; una mostra, l'altra trova un modo diverso di eseguirlo. La musica permette alla conoscenza di passare attraverso il gesto, l'ascolto, la ripetizione e la variazione. La piattaforma di studio conserva materiali e percorsi; l'incontro li rimette in movimento.

\section{Le decisioni che seguono da questa idea}
La coppia di postazioni è l'unità minima per valutare la promessa relazionale. Il collaudo di una sola non basta. Volto e mani devono essere leggibili e l'audio deve permettere di distinguere il proprio gesto da quello remoto. Il collegamento deve essere reciproco e riconoscibile, con la libertà di chiudere camera, microfono e sessione. Una registrazione disponibile per lo studio va dichiarata come registrazione: non può occupare furtivamente il posto di una persona.

La musica simultanea richiede una verifica specifica della tratta e dell'attività. La distanza non viene cancellata da una scelta grafica o da un algoritmo di predizione. Può però essere compresa, misurata e abitata attraverso diverse pratiche musicali. Un sistema che rende esplicite le proprie condizioni permette ai partecipanti di scegliere; uno che promette assenza di ritardo e restituisce un'esperienza incoerente consuma la fiducia su cui l'incontro si regge.

Lo studio individuale è una funzione piena. Deve poter cominciare senza un interlocutore disponibile e mantenere continuità con gli incontri: una frase da riprendere, una difficoltà da esplorare, una composizione da offrire alla visita successiva. Una piattaforma propria ha senso se rende questi passaggi più semplici; le risorse esterne hanno valore quando sono accessibili, pertinenti e utilizzabili nelle condizioni reali del luogo. L'AI resta un supporto eventuale e distinto dalla presenza umana.

La città entra nella relazione con i propri tempi e responsabilità. Occorrono orari compatibili, persone che preparino l'incontro, cura dello strumento e attenzione a chi attraversa il giardino senza voler partecipare. Collegare due città richiede che entrambe possano dare e ricevere: repertori, domande, scoperte e tempo. Una rete cresce bene se ciascun luogo può riconoscersi nel progetto comune.

\section{Che cosa dobbiamo poter apprendere dalla prima prova}
La prima coppia funzionante deve permettere di osservare se la presenza dell'altro è comprensibile e desiderabile; se la persona riesce a orientare camera, ascolto e attività senza perdersi nell'interfaccia; quali pratiche musicali sopportano le condizioni della tratta; che cosa resta appreso quando il partner e i suggerimenti non ci sono più. Queste domande richiedono misure diverse. Sentirsi vicini non equivale a tenere lo stesso tempo; tenere il tempo non dimostra automaticamente un apprendimento duraturo.

Un risultato può suggerire un cambio d'inquadratura, un altro percorso di rete, una diversa attività, più facilitazione o un ambiente di studio più semplice. Può anche mostrare che una normale configurazione di videolezione offre risultati equivalenti a minor costo. Il Portale deve rendere discutibile il proprio valore aggiunto: la forma, il rapporto con lo spazio pubblico e la continuità fra studio e incontro devono giustificare il lavoro necessario a realizzarli.

\section{Una forma di vicinanza da costruire}
Vorrei conservare, alla fine del progetto tecnico, la scena che lo ha generato: qualcuno alza gli occhi dalla tastiera e trova una persona che lo sta ascoltando. Non sa ancora che cosa suoneranno, né se riusciranno subito a capirsi. Sa che dall'altra parte c'è spazio per una risposta.

Il carattere poetico del Portale vive in questa apertura. Una città può offrirla a chi non dispone di uno strumento, di un insegnante vicino o di un compagno con cui provare. Può anche offrirla a un musicista esperto che desideri incontrare un altro modo di ascoltare. La tecnologia aumenta la portata dell'incontro; l'apprendimento cresce nella qualità dell'attenzione che le persone riescono a dedicarsi.

Questa è la promessa che vale la pena costruire con precisione: dare un luogo a una vicinanza possibile, lasciando che la musica attraversi la distanza e che ciascuno possa tornare al proprio studio con qualcosa ricevuto da un altro.
